\documentclass[tikz,border=6pt]{standalone}
\usepackage{ctex}
\usepackage{amsmath}
\usepackage{amssymb}
\usepackage{pgfplots}
\pgfplotsset{compat=1.18}
\usetikzlibrary{calc,angles,quotes,arrows.meta,positioning,decorations.markings}
\begin{document}
\begin{tikzpicture}[x=1.1cm]
% 对数刻度尺：位置 = lg(value)，跨度 1..100
% 取 lg：1->0, 2->0.301, 3->0.477, 4->0.602, 5->0.699, 10->1, 20->1.301, 100->2
% 用 8 倍放大把 lg 值映射成横坐标
\def\S{4.6} % 每个数量级的物理长度(cm)

% ===== 标尺主体 =====
\draw[very thick] (0,0) -- ({2*\S},0);
% 刻度：列出 value 与其 lg
\foreach \v/\lg in {1/0,2/0.30103,3/0.47712,4/0.60206,5/0.69897,6/0.77815,7/0.84510,8/0.90309,9/0.95424,10/1,20/1.30103,30/1.47712,40/1.60206,50/1.69897,100/2}{
  \draw[thick] ({\lg*\S},0) -- ({\lg*\S},0.22);
  \node[font=\footnotesize,anchor=south] at ({\lg*\S},0.24) {$\v$};
}
\node[anchor=north west,font=\small\bfseries] at (0,-0.05) {对数刻度（位置 $=\lg(\text{数值})$）};

% ===== 用三段“长度”演示 lg4 + lg2 = lg8 =====
% lg4=0.60206, lg2=0.30103, lg8=0.90309
\def\lgFour{0.60206}
\def\lgTwo{0.30103}
\def\lgEight{0.90309}

% 第一段：从 1 到 4，长度 = lg4
\draw[blue,very thick,{|[blue]}-{Stealth[blue]}] (0,-0.95) -- ({\lgFour*\S},-0.95)
   node[midway,below=1pt,font=\scriptsize,blue]{$\lg 4\approx0.602$};
% 第二段：接续 lg2，从 4 的位置再走 lg2，终点应到 lg8
\draw[red,very thick,{|[red]}-{Stealth[red]}] ({\lgFour*\S},-0.95) -- ({\lgEight*\S},-0.95)
   node[midway,below=1pt,font=\scriptsize,red]{$\lg 2\approx0.301$};

% 标出起点 1、中点 4、终点 8 的竖直对位虚线
\draw[blue!50,dashed] ({\lgFour*\S},0) -- ({\lgFour*\S},-0.95);
\draw[red!55,dashed] ({\lgEight*\S},0) -- ({\lgEight*\S},-0.95);
\draw[black!45,dashed] (0,0) -- (0,-0.95);

% 终点强调
\node[anchor=north,font=\scriptsize,blue] at (0,-1.35) {$1$};
\node[anchor=north,font=\scriptsize] at ({\lgFour*\S},-1.35) {$4$};
\node[anchor=north,font=\scriptsize,red] at ({\lgEight*\S},-1.35) {$4\times2=8$};

% 合计长度大括号
\draw[decorate,decoration={brace,amplitude=6pt,mirror},thick]
   (0,-1.75) -- ({\lgEight*\S},-1.75);
\node[anchor=north,font=\scriptsize] at ({\lgEight*\S/2},-1.95)
   {两段长度相加 $=\lg 4+\lg 2=\lg 8\approx0.903$};

% ===== 结论注释框 =====
\node[fill=white,fill opacity=0.85,text opacity=1,draw=black!55,rounded corners,
      font=\small,align=center,inner sep=5pt] at ({1.5*\S},1.55)
  {计算尺原理：\\$\lg(x\,y)=\lg x+\lg y$\\\emph{乘法} 在对数刻度上变成 \emph{长度相加}};

\end{tikzpicture}
\end{document}
